\chapter{Building atomic structures}\label{ch:builders}
\section{Common builder workflow}
Choose the appropriate command under \ui{Build}, supply a periodic reference where required, set parameters, inspect the preview, and run the build action. Use Load reference in Nanocrystal, CSL, or Polycrystal; Load structure or model in Custom Structure; and Add structure in Merge. These buttons open the shared file picker, while file drops remain available. Read status messages before saving the result. Source previews and generated previews serve different purposes: changing a parameter may update an illustrative preview while the atomic result still needs rebuilding.

Successful Build-menu results open in a new structure tab through the main application's builder callback, preserving the source tab. Merge uses this same new-tab route. Select the generated tab before saving or running another tool. Direct atom edits, Cell Sculptor, Dislocation replacement, and Interstitial insertion update the current structure instead. In the detailed controls below, ``updates the main scene'' refers to the displayed result; it does not mean that every builder overwrites the original tab.

Use modest cells first. Replicating a structure multiplies its atom count; a $10\times10\times10$ replication of a four-atom cell contains 4,000 atoms before cuts. Increasing all three dimensions by two multiplies the workload by eight. Once the geometry is correct, increase size and inspect boundary contacts. A builder does not automatically relax strained bonds or remove every physically undesirable local environment.

\section{Bulk Crystal}
\uishot{bulk-dialog.png}{450bp 175bp 450bp 192bp}{The Bulk Crystal dialog at its initial settings. Change the default H site to Cu and the default lattice parameter to 3.61 Angstrom for the FCC Cu walkthrough. The element button opens a periodic-table picker.}
\textbf{Purpose.} Generate a symmetry-expanded periodic crystal from lattice parameters and asymmetric-unit sites. Inputs are the crystal system, space group, cell dimensions/angles, and element/fractional-coordinate rows.

\textbf{Procedure.} Select one of the seven crystal systems and a compatible space group (1--230). Set independent lattice parameters; check constraints implied by the selected system. Add each crystallographically independent site using an element and direct coordinates. Use the row deletion or clear controls to remove unwanted defaults. Build, inspect the expanded sites, and save a periodic format.

Fractional coordinates are coefficients in $\mathbf r=f_a\mathbf a+f_b\mathbf b+f_c\mathbf c$. They are not distances in \angstrom. Equivalent positions are generated by symmetry; entering a full conventional basis as though every site were independent can duplicate sites. Inspect the final count and short contacts.

\textbf{Worked examples.} FCC Cu with group 225, $a=3.61$ \angstrom, and Cu at $(0,0,0)$ gives four atoms. BCC Fe with group 229, $a=2.87$ \angstrom, and Fe at $(0,0,0)$ gives two. These examples exercise symmetry expansion; they do not determine a thermodynamically optimal lattice parameter.
\begin{lstlisting}
AtomForge --build bulk --system cubic --spacegroup 225 --a 3.61 --atom "Cu 0 0 0" --output cu_fcc.cif
AtomForge --build bulk --system cubic --spacegroup 229 --a 2.87 --atom "Fe 0 0 0" --output fe_bcc.cif
\end{lstlisting}
For lower-symmetry cells, set \code{--b}, \code{--c}, and angles \code{--alpha}, \code{--beta}, \code{--gamma}. Repeat \code{--atom} for independent sites. See \code{--help} for accepted system names.

\section{Substitutional Solid Solution}
\textbf{Purpose.} Assign species to existing lattice sites with requested fractions. This changes chemical labels on a host structure; it does not search for a special quasirandom structure or optimise chemical short-range order.

\textbf{Procedure.} Open a host with enough sites to represent the composition, then choose \ui{Build > Substitutional Solid Solution}. Select the source and species percentages. The GUI uses at\%, with the final row automatically balancing the total to 100\%; the CLI uses fractions summing to one. Choose a fixed seed for reproducibility. Generate the alloy and check the resulting species counts before export. For small hosts, integer atom counts limit the achievable composition.

The example host is a $5\times5\times5$ repetition of conventional FCC Cu: 500 sites. The 70\% Cu / 30\% Ni assignment uses seed 42. Keep the same host atom ordering and seed when reproducing an assignment; changing the host can change which particular sites are substituted.
\begin{lstlisting}
AtomForge --build sss --input cu_host.cif --frac "Cu=0.7,Ni=0.3" --seed 42 --output cu_ni_alloy.cif
\end{lstlisting}
The menu is controlled by \code{ATOMFORGE_ENABLE_SSS_BUILDER}, enabled by default. A seed of zero in the CLI requests a time-dependent seed. After substitution, an external relaxation may be needed to account for atomic size mismatch.

\section{CSL Grain Boundary}
\textbf{Purpose.} Construct a coincidence-site-lattice bicrystal from a compatible cubic periodic reference. A coincidence index $\Sigma$ describes lattice coincidence, not a calculated grain-boundary energy.

\textbf{Procedure.} Load the reference cell in the grain-boundary dialog. Enter a rotation axis $[u\ v\ w]$ and select a supported $\Sigma$ and candidate boundary plane. Set the thickness/repetition of each grain, any gap and vacuum, and an overlap-removal distance. Build and inspect both the interface and periodic cell boundaries. Save the bicrystal and relax it externally before comparing energies.

\textbf{Controls.} The axis is a crystallographic direction. The plane selector chooses from generated candidates. The CLI's \code{--plane} is a zero-based candidate index, not a Miller index. \code{--uca} and \code{--ucb} control stacking repetitions for the two grains. \code{--gap} separates the grains; \code{--vacuum} adds vacuum on each side. \code{--overlap} removes close contacts using a distance in \angstrom; an excessive cutoff can delete desirable interfacial atoms.

\begin{lstlisting}
AtomForge --build gb --input cu_fcc.cif --axis "0 0 1" --sigma 5 --plane 0 --uca 3 --ucb 3 --overlap 1.5 --output cu_sigma5.cif
\end{lstlisting}
This manual's generated example contains 60 atoms. The small thickness makes it a geometry demonstration, not a converged boundary-energy supercell. If \code{--sigma} is omitted, candidate search uses \code{--sigmamax}; specify the exact value when reproducibility matters. \code{--conventional} skips the primitive-cell conversion. Check both lattice compatibility and the selected candidate before interpreting a failed build as a file-format problem.

\section{Nanocrystal: shape cuts}
\textbf{Purpose.} Fill a finite shape with atoms from a periodic crystal. Load a periodic reference, select \ui{Shape Cut}, choose a shape, set dimensions and centre, and use automatic replication unless you need explicit repetition counts. Build the nanocrystal and inspect the output bounding cell.

\begin{longtable}{p{.29\linewidth}p{.62\linewidth}}
\toprule Shape & Main dimensions \\ \midrule\endhead
Sphere & Radius. \\
Ellipsoid & Radii along the three local axes. \\
Box & Half-lengths along the three axes; full side lengths are twice these values. \\
Cylinder & Radius, full height, and cylinder axis. \\
Octahedron & Size parameter for the octahedral cut. \\
Truncated octahedron & Size and truncation parameters. \\
Cuboctahedron & Size parameter for the cuboctahedral cut. \\
\bottomrule
\end{longtable}
\ui{Auto-center from atoms} derives a centre from the reference. Manual centring changes which lattice sites intersect the boundary and can change surface termination and atom count. \ui{Auto-replicate} chooses enough source repetitions for the cut; manual repetitions that are too small truncate the intended object. \ui{Rectangular cell} and vacuum settings provide an enclosing output cell; this does not make a nanoparticle a bulk material.

\begin{lstlisting}
AtomForge --build nano --input cu_fcc.cif --shape sphere --radius 10 --vacuum 5 --output cu_sphere.xyz
\end{lstlisting}
The example contains 336 atoms. Use \code{--rx/--ry/--rz} for ellipsoid radii, \code{--hx/--hy/--hz} for box half-lengths, and \code{--cylradius/--cylheight/--cylaxis} for the cylinder. Explicit source replication uses \code{--repa}, \code{--repb}, and \code{--repc}; zero selects automatic sizing in the CLI.

\section{Nanocrystal: Wulff construction}
Select \ui{Wulff Construction} in the nanocrystal dialog. Add facet families, each with nonzero $(hkl)$ indices and a positive relative surface energy. AtomForge expands families using the reference symmetry, previews the bounding planes, and cuts the periodic crystal by the resulting shape.

\begin{enumerate}
\item Load the periodic reference and check its orientation and symmetry.
\item Add, for example, $(100)$ and $(111)$ families with illustrative energies 1.0 and 0.9. These are demonstration ratios, not material data.
\item Set \ui{Max radius (A)}. The implementation scales plane distances by energy ratios, placing the highest-energy family at this maximum distance. This control is not necessarily the circumscribed vertex radius of the final polyhedron.
\item Inspect the \ui{Wulff Plane View}; family colours help identify the contribution of each family. Remove redundant families or add missing ones as required.
\item Build the atomic nanocrystal. Inspect discrete surface termination separately from the continuous plane preview, then export.
\end{enumerate}
The construction requires a bounded intersection of planes. A family may disappear from the final shape when other planes enclose it. Energy ratios must come from comparable surface calculations if a physical equilibrium-shape interpretation is intended. Atomic discretisation can break an ideal continuous shape's apparent symmetry in a very small particle. The same construction is available through Python \code{af.wulff} and native \code{--build nano --shape wulff}; use repeated \code{--facet "h k l energy"} arguments for the families.

\section{Custom Structure from OBJ/STL}
\textbf{Purpose.} Fill a closed mesh with lattice sites. The mesh is a spatial boundary, not a collection of atoms. Use a watertight OBJ or STL, a periodic atomic reference, a physical scale, and an orientation.

\textbf{Procedure.} Open \ui{Build > Custom Structure}. Load/drop the reference and mesh into their corresponding inputs. Set the mesh scale in \angstrom\ per model unit, then adjust rotation or the Miller-direction alignment. Select enough replication to cover the mesh, set vacuum, and build. Inspect thin features and the atom count; a feature narrower than the lattice spacing may contain no sites.

\begin{lstlisting}
AtomForge --build custom --input cu_fcc.cif --mesh octahedron.obj --scale 12 --vacuum 5 --output cu_mesh.xyz
\end{lstlisting}
The included octahedron has six vertices at the positive and negative unit axes and eight triangular faces. At scale 12 it produces 224 Cu atoms in this example. No Blender installation is required. To use Blender, create a closed mesh, apply the intended object transforms before export, export OBJ/STL, and verify the scale in AtomForge. Inspect for holes, self-intersections, and reversed or inconsistent faces when the fill is missing or unexpected.

The CLI provides \code{--rotx}, \code{--roty}, and \code{--rotz} in degrees. \code{--miller "h k l"} aligns a crystallographic direction with $+Z$ and overrides Euler orientation. A crystallographic direction $[hkl]$ and a plane normal $(hkl)$ differ in a general non-cubic cell; choose the alignment convention deliberately.

\section{Polycrystal}
\textbf{Purpose.} Populate a box with grains cut by a Voronoi partition. Inputs include a periodic reference, box dimensions, grain count, random seed, and optional orientations. The construction creates geometric grain boundaries; it does not optimise their energies.

\textbf{Procedure.} Load a reference, enter box lengths and grain count, set the seed, and choose orientations or random assignment. Generate, inspect the grain sizes and contacts, and save. Use orientation colouring when the output retains orientation metadata. Reopening a format that lacks grain metadata can preserve atoms while losing those colours.

\begin{lstlisting}
AtomForge --build poly --input cu_fcc.cif --sizex 25 --sizey 25 --sizez 25 --grains 4 --seed 7 --output cu_polycrystal.cif
\end{lstlisting}
This example contains 1,329 atoms. The CLI accepts repeated \code{--euler "phi1 Phi phi2"} triplets in Bunge convention, one per grain; unspecified orientations are random. Use the same seed, dimensions, reference, and ordering to reproduce a sample. Increase the box relative to grain count when grains become too small to have a meaningful crystalline interior.

\section{Amorphous Structure}
\textbf{Purpose.} Produce a random packing subject to distance constraints. Specify elements and counts, a mass density or explicit box, pair-distance constraints, and a seed. This is an initial configuration generator; random packing alone does not establish a realistic glass network.

\textbf{Procedure.} Add the required species/count rows. Set density in $\mathrm{g\,cm^{-3}}$ or explicit box dimensions. Check the minimum separation settings; overly strict cutoffs at a high density can make placement impossible. Set a fixed seed, generate, and inspect the reported inserted count and any placement failure. Relax or equilibrate the configuration with a suitable external method before structural or energetic claims.

\begin{lstlisting}
AtomForge --build amorphous --element "Si 40" --element "O 80" --density 1.5 --seed 7 --output sio2_pack.xyz
\end{lstlisting}
The example has 120 atoms with a 1:2 Si:O ratio. The deliberately low packing density is useful for demonstrating construction and is not a claim about silica's equilibrium density. \code{--boxa/--boxb/--boxc} override density-derived dimensions. \code{--mindist "Z1 Z2 distance"} adds a pair constraint using atomic numbers; \code{--covtol} controls the covalent-radius fallback; \code{--attempts} bounds placement attempts. If placement fails, examine the geometry and constraints before merely raising the attempt limit.

\section{Optional and unavailable builder menus}
The stacking-fault builder is compiled into the desktop menu only with \code{ATOMFORGE_ENABLE_SFE_BUILDER=ON}; the default is OFF. Its purpose is to construct faulted configurations, not to calculate a stacking-fault energy. In an enabled build, open \ui{Build > Stacking Faults}, supply a reference with \ui{Use Current Scene} or the source workflow, and click \ui{Detect Structure}. Select a supported Plane, choose \ui{Smallest Unit Cell} or \ui{Orthogonal Cell}, and set Layers (2--96). Click \ui{Generate Sequence}, use the Structure slider to inspect its members, and choose \ui{Use Selected Structure} to adopt one. \ui{Export All Structures} writes the sequence to the specified Folder/Subfolder; \ui{Use Source Folder} supplies a convenient base when available. Validate the stacking geometry and calculate energies externally. This optional path was inspected in source but not enabled in the screenshot executable.

The interface-builder dialog has a drawing call but no current main-menu entry to open it. Its source/API capabilities are described in Appendix~\ref{ch:coverage}. Use the exposed merge/transform workflows for manual assembly in the current executable.

\section{Worked study: compare bulk and finite Cu environments}
This exercise connects building, file inspection, and structural analysis. It uses only generated structures and does not require an electronic calculation.
\begin{enumerate}
\item Open \code{examples/cu_fcc.cif}. Inspect its four sites and cubic cell. The cell edge is 3.61 \angstrom. Save a working copy before changing it.
\item Open the supplied \code{cu_host.cif}, or create it by repeating the conventional cell $5\times5\times5$ through Python. It contains 500 atoms. Use this larger periodic host for bulk neighbour statistics.
\item Open \code{cu_sphere.xyz} separately. It contains 336 atoms from the same reference lattice, cut at radius 10 \angstrom. Rotate it and inspect surface sites. The smaller count is caused by the finite shape, not by a different crystal symmetry.
\item Run RDF on the periodic host with Cu as both species, PBC enabled, and no smoothing. Check the first-neighbour peak near 2.553 \angstrom. Begin with a cutoff that is comfortably smaller than the box dimensions.
\item Run RDF on the finite sphere with PBC disabled. Compare raw counts and cumulative coordination before comparing normalised curves: the particle has surface atoms with incomplete neighbour shells and its enclosing vacuum does not define a bulk atomic density.
\item Open \code{cu_ni_alloy.cif}. Verify 350 Cu and 150 Ni atoms, then inspect Cu--Cu and Cu--Ni partial distributions and the SRO table. The unrelaxed substitution changes labels while retaining ideal host positions; it cannot demonstrate size-mismatch relaxation.
\end{enumerate}
Keep the source filenames, cutoffs, species choices, PBC setting, bin count, and smoothing with each plot. If the bulk nearest-neighbour peak moves unexpectedly, first check lattice scaling and units. If coordination falls at the surface of the sphere, inspect its geometric boundary before treating the change as disorder. If the alloy has nonzero measured SRO, compare seeds and finite-size effects before interpreting it as evidence of an energetic ordering preference.

This comparison illustrates why a common visual appearance does not guarantee comparable statistics. The bulk, particle, and alloy answer different questions even though they originate from the same Cu cell. Choose analysis settings according to the geometry and composition of each saved structure.
